\subsection{Error patterns by recent rainfall}
The final 2024 holdout contains only stations seen during training; all
6,374,476 prediction rows have \texttt{unseen\_station=false}. It therefore
does not measure transfer to new stations. To expose performance differences
within the observed network, Table~\ref{tab:onset} separates samples with no
rain in the rolling 30-minute history from those with any rain in that window.
The flag is derived from \texttt{sum\_30}, which includes the current reading
and the prior five readings. Metrics use final PAVA scores and the same fixed
0.5 probability cutoff.

\begin{table}[htbp]
\centering\scriptsize
\caption{Final PAVA classification by recent-rainfall regime in held-out
2024. The rolling 30-minute history includes the current reading and previous
five readings. Precision and F1 are undefined when no positive predictions
were made.}
\label{tab:onset}
\begin{tabular}{r l r r r r r}
\toprule
Threshold (mm) & Recent history & Windows & Positives & Prevalence (\%) &
Precision & Recall / F1 \\
\midrule
0.5 & No rain & 5,969,224 & 69,263 & 1.16 & --- & 0.000 / --- \\
0.5 & Any rain & 405,252 & 161,808 & 39.93 & 0.7417 & 0.7524 / 0.7470 \\
10.0 & No rain & 5,969,224 & 6,894 & 0.12 & --- & 0.000 / --- \\
10.0 & Any rain & 405,252 & 17,436 & 4.30 & 0.5859 & 0.2690 / 0.3687 \\
\bottomrule
\end{tabular}
\end{table}

At both thresholds, the model made no positive predictions after a dry
30-minute history, giving zero recall in that subgroup. At 0.5~mm, recall was
0.7524 when there had been recent rainfall. The fixed 0.5 cutoff therefore
behaves very differently in dry and recently wet conditions.

